% Reference entries, R (AN-11).  Law 1: os/iword/app/face_a.c
% retrieve_named (a missing file, a file that is not a WordPerfect 5.1
% document, a file larger than 64K; onto a clear screen the file becomes the
% document, into a document its text goes in at the cursor), os/iword/fs/
% wp5.c iw_wp5_read (only codes the corpus verified are accepted),
% spec/iword/wp5_format.h (HRt 0x0A, SRt 0x0D, DSRt 0x90), edit.c (wrap,
% the DSRt of Justification), codes.c iw_reveal_pane and face_a.c
% paint_reveal (the Reveal Codes screen: 10 document rows, the status line,
% the ruler, the 10-row codes area, the restore message; spec/iword/
% reveal_codes.h).  Shots: ../wordperfect51-decomp goldens MINT01..05,
% mint MINT06..32.
% GATES: test-pubs-iword test-iword-face test-iword-wp5-roundtrip test-iword-reveal-codes test-flair-iword test-flair-iword-codes test-pubs-flair-screens
\Gates{test-pubs-iword test-iword-face test-iword-wp5-roundtrip test-iword-reveal-codes test-flair-iword test-flair-iword-codes test-pubs-flair-screens}
\begin{Feature}{Retrieve}
Retrieve brings a document you have saved back onto the screen, so that you
can read it or change it.  You can also
retrieve a document into the one on the screen, to combine the two.

To retrieve a document,
\begin{WPSteps}
\item Clear the screen if it holds a document you do not want to combine
with the one you retrieve (see \InRef{Exit}).
\item Press \FK{Retrieve}{Shift-F10}.  The status line shows
\Scr{Document to be Retrieved:}.
\PullDown{Select \Mi{R}etrieve from the \Mi{F}ile menu.}
\item Type the name of the document and press \K{Enter}.
\end{WPSteps}
\Documents{wpkey}{Shift-F10}
On a clear screen the document you retrieve becomes the document you are
working on: its name is shown on the status line and the cursor is at its
top.

\Notes
\Topic{Into Another Document}
If the screen already holds text, the retrieved document's text is put into
it at the cursor, and the document on the screen keeps its own name.  The
two documents are now one; the file you retrieved is not changed.

\Topic{Messages}
\Scr{ERROR: File Not Found} means there is no document by that name on the
data volume: check the name and type it again.
\Scr{ERROR: Incompatible File Format} means the file is not a WordPerfect
5.1 document that InitechWord can read.
\Scr{ERROR: Document too large} means the file is larger than InitechWord
can hold (see Appendix H).  Press any key to return to the prompt.

\Topic{WordPerfect 5.1 Documents}
InitechWord's documents are WordPerfect 5.1 documents, and InitechWord
retrieves documents written by WordPerfect 5.1 itself, keeping every code
it holds (see \InRef{Codes}).  A WordPerfect document that holds a code
InitechWord does not know is refused with \Scr{ERROR: Incompatible File
Format}, rather than changed.

\SeeAlso{Exit; Filenames; Save}
\end{Feature}

\begin{Feature}{Return, Deletable Soft}
A deletable soft return ends a line without ending the paragraph.
InitechWord puts one in when you set the justification with the cursor in
the middle of a line, so that the text after the cursor starts a new line
in the new justification (see \InRef{Justification}).

Reveal Codes shows it as \Code{[DSRt]}.  To join the two lines again, delete
the code: InitechWord deletes it without asking.

\SeeAlso{Justification; Return, Hard; Return, Soft}
\end{Feature}

\begin{Feature}{Return, Hard}
A hard return ends a line where you choose, and ends the paragraph.  You put
one in by pressing \K{Enter}.  Use it at the end of a paragraph, after each
line of an address or a list, and to make blank lines.

To insert a hard return,
\begin{WPSteps}
\item Press \K{Enter}.
\end{WPSteps}
\Documents{wpkey}{Enter}

\Notes
\Topic{Code}
A hard return is the code \Code{[HRt]}.  Delete it to join a line to the
line after it (see \InRef{Delete}).

\Topic{Shift-Enter and Alt-Enter}
Shift-Enter and Alt-Enter put in a hard return as Enter does.  Ctrl-Enter
puts in a hard page (see \InRef{Page Break, Hard}).

\SeeAlso{Page Break, Hard; Return, Soft; Word Wrap}
\end{Feature}

\begin{Feature}{Return, Soft}
A soft return is the end of a line that InitechWord makes for you when a
line is full (see \InRef{Word Wrap}).  It takes the place of the space
between the last word that fits on the line and the word that goes to the
next line.  Reveal Codes shows it as \Code{[SRt]}.

You never type a soft return and you need not delete one: whenever you
change a paragraph, InitechWord wraps it again and moves its soft returns to
where the lines now end.

\SeeAlso{Return, Deletable Soft; Return, Hard; Word Wrap}
\end{Feature}

\begin{Feature}{Reveal Codes}
Reveal Codes shows the hidden codes of the document.  The screen divides:
the upper part shows the document as usual, and the lower part shows the
same lines with every code in square brackets.  Use it to see why text
looks as it does, and to delete a code.

To reveal codes,
\begin{WPSteps}
\item Press \FK{Reveal Codes}{Alt-F3}.
\item Move the cursor as usual.  It moves in both parts of the screen; in
the codes area it is the highlighted character or code.
\item Press \FK{Reveal Codes}{Alt-F3} again to return to the full editing
screen.
\end{WPSteps}
\Documents{wpkey}{Alt-F3}
\FlairShot{reveal}{640}{440}{Reveal Codes: the bold codes around a word}

The Reveal Codes screen has five parts, from the top:
\begin{description}
\item[The document] ten lines of the document, as on the editing screen.
\item[The status line] as on the editing screen.
\item[The ruler] the margins and tab stops in effect at the cursor.
\item[The codes area] ten rows: the text of the lines around the cursor,
starting three lines above it, with the codes shown.  A line longer than the
screen runs on into the next row.
\item[The message] \Scr{Press Reveal Codes to restore screen}.
\end{description}

\Notes
\Topic{Deleting Codes}
In Reveal Codes, Backspace and Delete delete codes without asking, as they
delete characters.  Deleting either code of an appearance or size deletes
its partner too.  See \InRef{Codes}.

\Topic{Editing in Reveal Codes}
Everything you can do on the editing screen you can do in Reveal Codes:
type, block, move, format.  Both parts of the screen change as you work.

\Topic{The Cursor in the Codes Area}
The highlight covers the character or the whole code the cursor is on.  At
the very end of the document it is one empty cell after the last code.

\Topic{The Ruler}
The left margin is shown by \Scr{\{} and the right margin by \Scr{\}}; a
margin that does not fall on a tab stop is shown by \Scr{[} or \Scr{]}
instead.  The small triangles are the tab stops, every half inch from the
left margin.

\SeeAlso{Codes; Margins, Left and Right; Tab}
\end{Feature}
